\documentclass[a4paper]{article}

\usepackage[a4paper, top=8mm, bottom=8mm, left=8mm, right=8mm]{geometry}

\usepackage{polyglossia}
\setdefaultlanguage[babelshorthands=true]{russian}
\setotherlanguage{english}

\usepackage{fontspec}
\setmainfont{FreeSerif}
\newfontfamily{\russianfonttt}[Scale=0.7]{DejaVuSansMono}

\usepackage[tiny, compact]{titlesec}

\usepackage{titling}
\setlength{\droptitle}{-1cm}
\pretitle{\begin{center}\begin{bfseries}\Large}
\posttitle{\par\end{bfseries}\end{center}}
\preauthor{\begin{center}\normalsize}
\postauthor{\par\end{center}\vspace{-1.8cm}}

\usepackage{hyperref}
\usepackage{bookmark}
\usepackage{csquotes}
% Сюда можно дописывать нужные вам пакеты.

\title{Менеджер метаданных для системы MarLine}

% Уточните отчество.
\author{Потапов Денис}

% Дата много места занимает, её лучше не указывать.
\date{}

\begin{document}

\maketitle

\begin{flushright}
    Группа: \emph{25.Б71-мм}

    Кафедра: \emph{Кафедра системного программирования СПбГУ}

    % Степень/звание и должность научника мы лучше вас знаем, их не указываем.
    Научный руководитель: \emph{Гориховский Вячеслав Игоревич}

    % Укажите официальное место работы и должность по трудовой.
    Консультант: \emph{Дмитриевцев Алексей Сергеевич}

    Младший инженер по разработке ПО ООО "КНС ГРУПП"

    Номер семестра практики: \emph{3}
\end{flushright}


\section{Постановка задачи}

Работа выполняется в рамках open source-проекта MarLine --- системы хранения
данных с двухуровневой дедупликацией. На первом уровне контентно-ориентированное
разбиение (CDC) задаёт границы блоков так, чтобы они оставались устойчивыми к
сдвигам данных при изменении файлов, а на втором уровне Similarity-Based Chunking
подбирает для каждого блока похожий базовый блок и хранит его в виде дельты.

Одним из методов построения такой дельты является Palantir --- наследник подхода
Odess, основанный на иерархии тиров (уровней) и суперфич: многоуровневая структура
позволяет группировать блоки с разной степенью схожести и последовательно сужать
множество кандидатов поиска пары для дельта-кодирования.

Ключевую роль в этом механизме играет менеджер метаданных --- компонент,
отвечающий за хранение и поиск метаданных блоков в подсистеме Palantir. Он
обеспечивает поиск как точно совпавших блоков, так и потенциально схожих, а также
управляет жизненным циклом записей. Остальным конвейерам менеджер предоставляет
единый API: добавление блока, поиск по отпечатку и по суперфичам, завершение
версии.

\textbf{Целью работы является} разработка и внедрение менеджера метаданных в
алгоритм Palantir.

Для достижения цели требуется решить следующие задачи:

\begin{enumerate}
    \item Изучить идеи реализации Palantir в оригинальной статье и спроектировать
          архитектуру менеджера метаданных.
    \item Интегрировать инвертированный индекс в менеджер метаданных.
    \item Реализовать разбиение потока байт на поколения бэкапов.
    \item Интегрировать менеджер в алгоритм Palantir.
    \item Оформить отчёт, документацию и pull request в репозиторий проекта.
\end{enumerate}

Без менеджера метаданных поиск базового блока в Palantir не масштабируется,
поэтому результаты работы необходимы для работоспособности подсистемы.
В выполнении работы заинтересована команда разработки проекта MarLine.
\section{Обзор}

\section{Описание предлагаемого решения}

% Раздел пишется в свободной форме, с минимумом технических подробностей (можно приводить ссылки на более подробные документы).

%Тут надо кратко описать идею решения, спозиционировать его относительно предыдущих наработок\footnote{Если на них надо сослаться, то в подстраничной сноске, URL: \url{http://example.com} (дата обращения: не забываем указывать).} и существующих аналогов (тоже предельно кратко), по возможности обосновать принятые решения.

% Вот это важно, опишите кратко, что использовали
%Реализация решения осуществлялась с использованием \dots (тут описать используемые технологии).

\section{Апробация / Эксперименты}

% Название секции можно поменять на более подходящее вашей работе

%Название этой секции --- либо \enquote{Апробация}, либо \enquote{эксперименты} (если содержательно есть и то и то, можно две отдельные секции сделать). Тут опишите, как проводилась проверка того, что получилось что-то адекватное. Если есть отзывы пользователей, отлично (особенно если они в структурированном виде, типа анкеты SUS). Если есть отзывы команды, консультанта и т.п., хорошо. Если есть только тесты. то хотя бы их тут опишите.

%Если работа подразумевала замеры чего-то (например, производительности), приведите тут итоговые результаты и выводы, и дайте ссылку на таблицы с данными. Не забывайте про матстат (приводите матожидание и среднеквадратичное отклонение, не просто результат замера).

%Приведите выводы из эксперимента.

\section{Заключение}

%В ходе выполнения работы получены следующие результаты:

% ниже списком перечисляете то, что выносится как результат практики.

%\begin{itemize}
%    \item \dots
%    \item \dots
%    \item \dots
%\end{itemize}

%Материалы, иллюстрирующие выполнение работы, размещены по адресу: \url{http://www.example.com} (если есть и применимо --- скриншоты, видео на 1-2 минуты и т.п. с демонстрацией; может быть очень полезно даже для консольных утилит или библиотек).

%Репозиторий/ссылка на пуллреквест: \url{http://www.example.com} (это обязательно; ссылок может быть несколько).

%Техническая документация: \url{http://www.example.com} (если применимо --- для пуллреквестов можно техническую сторону прямо в комментарии к пуллреквесту указать, но можно и в отдельную вики-страницу вынести).

%Решение внедрено в производственные процессы компании/прошло процесс ревью/находится на ревью/находится в разработке и т.п. В общем, кратко сформулируйте, что в итоге получилось и в каком статусе результат, чтобы не было ощущения, что оно сделано, положено на GitHub и забыто.

\end{document}
